
\section{9. Discussion}

\subsection{9.1 Key Takeaways}

\subsubsection{9.1.1 Takeaway 1: Learned Spans are Interpretable}

The learned spans reveal a \textbf{structural} pattern in the model: some heads 
are naturally local, others global. This is not visible from fixed-window 
methods.

\subsubsection{9.1.2 Takeaway 2: Span Regularization Works}

The L1 penalty causes spans to shrink gracefully, with minimal accuracy loss 
at $\lambda = 0.01$.

\subsubsection{9.1.3 Takeaway 3: No Speedup Without Hard Truncation}

The soft ramp is differentiable but \textbf{does not speed up inference} because 
all scores are still computed. To realize speedup, one would:
1. Train with soft ramp.
2. Convert to hard window at inference.
3. Use a sparse kernel.

\subsection{9.2 When to Use ASA-Span}

\subsubsection{Recommended}

\begin{itemize}
\item \textbf{Interpretability studies} (head structure analysis)
\item \textbf{Base for hard-window pruning}
\item \textbf{Long-context fine-tuning}
\item \textbf{Combined with other methods}
\end{itemize}

\subsubsection{Not Recommended}

\begin{itemize}
\item \textbf{Immediate speedup needed} (use hard window directly)
\item \textbf{Parameter-constrained} (use A3/AASVD)
\item \textbf{Very short sequences} (no benefit)
\end{itemize}

\subsection{9.3 Choosing Max Span}

\begin{table}[H]
\centering
\small
\begin{tabular}{ll}
\toprule
\textbf{Sequence length} & \textbf{Recommended $S_\text{max}$} \\
\midrule
128 & 64 \\
512 & 128 \\
2048 & 256 \\
8192 & 512 \\
\bottomrule
\end{tabular}
\end{table}

\textbf{Rule of thumb.} $S_\text{max} = L / 4$ to $L / 2$.

\subsection{9.4 Choosing Ramp Size}

\begin{table}[H]
\centering
\small
\begin{tabular}{ll}
\toprule
\textbf{Use case} & \textbf{Recommended $R$} \\
\midrule
Aggressive pruning & 8 \\
Balanced & 16 \\
Smooth & 32 \\
Very smooth & 64 \\
\bottomrule
\end{tabular}
\end{table}

\subsection{9.5 Cost-Benefit}

For DistilBERT at $L = 128$:

\begin{table}[H]
\centering
\small
\begin{tabular}{llll}
\toprule
\textbf{Config} & \textbf{Accuracy} & \textbf{Mean span} & \textbf{Reduction} \\
\midrule
Baseline & 90.71\% & -- & -- \\
ASA-Span $\lambda=0$ & 90.5\% & 126 & 0\% \\
ASA-Span $\lambda=0.01$ & 89.5\% & 68 & 47\% \\
ASA-Span $\lambda=0.1$ & 86.0\% & 32 & 75\% \\
\bottomrule
\end{tabular}
\end{table}

\subsection{9.6 Conversion to Hard Window}

To realize speedup, one can:

1. \textbf{Train} with soft ramp.
2. \textbf{Extract} learned spans $\{z_h\}$.
3. \textbf{Round} to nearest power of 2 (for kernel alignment).
4. \textbf{Replace} soft ramp with hard band mask.
5. \textbf{Use sparse kernel} (e.g., block-sparse or flash attention with local window).

\subsection{9.7 Open Questions}

1. \textbf{Learned ramp size:} Can we learn $R$ per head?
2. \textbf{Interaction with content:} Does distance-based sparsity combine with 
   content-based (e.g., ASA-Sparse)?
3. \textbf{Cross-layer transfer:} Do spans transfer between tasks?
4. \textbf{Long-context scaling:} How do spans evolve at $L = 4096$?

\subsection{9.8 Broader Impact}

\subsubsection{Positive}
\begin{itemize}
\item \textbf{Interpretability:} Reveals head structure.
\item \textbf{Foundation for pruning:} Basis for hard-window deployment.
\item \textbf{Efficiency:} Enables fixed-window deployment.
\end{itemize}

\subsubsection{Negative}
\begin{itemize}
\item \textbf{Dual-use:} Enables surveillance at scale.
\item \textbf{Concentration:} Favors large labs.
\end{itemize}

\subsubsection{Environmental}
Total compute: ~12 GPU-hours, ~{{CO2}} kg CO$_2$eq.